\subsection{How Advent took shape}

A fourth-century Roman calendar records the birth of Christ at Bethlehem on
December 25. The entry is at the head of a list of anniversaries of martyrs,
set down by date in the \emph{Chronography of 354} (\emph{Depositio
Martirum}). When churches elsewhere begin to tell us how to spend the weeks
around the Nativity, they tell us about different things people were doing. In
380 the Council of Zaragoza bid people gather daily in church from December 17
through Epiphany, when some had been withdrawing to private houses, country
estates or the mountains (canon 4).

In fifth-century Rome St Leo the Great preached on a fast in December. By
then, as he reminds his congregation, the crops had been brought in. A
household with food put by had something to give thanks for and something to
share. Leo counsels each person to use his abundance by becoming “more
abstinent in himself and more open-handed towards the poor” (Sermon 16.2). In
another sermon he opens with the words Jesus speaks immediately after the
Gospel of the First Sunday: take heed lest excess and the cares of life weigh
down your hearts (Luke 21:34). The fast was one of four over the course of the
year (Sermon 19.1--2). The winter Ember Days of the Roman calendar continue
that seasonal observance. Their Masses in Advent, as they stand in the 1962
Missal, bring the fast very near to the events of the Nativity.

Gregory of Tours tells us that the schedule drawn up by his fifth-century
predecessor, Bishop Perpetuus, could still be read in his own day. Its entry
for the approach to Christmas runs, “From the burial of our lord Martin to the
Nativity of our Lord, three fasts each week” (\emph{Histories} 10.31). Martin
was the bishop whose tomb was in Tours. The observance began with his feast in
November and reached forward to Christmas. At Mâcon late in the sixth century
the bishops specified Monday, Wednesday and Friday for the fast and directed
that Mass be celebrated in the Lenten order during that time (First Council of
Mâcon, canon 9).

In the eighth century a Roman Mass book was sent to the Franks. In it, the
Gregorian Sacramentary, we find the two petitions priests still make on the
first two Sundays in the 1962 Missal. On the first the priest asks the Lord to
stir up his power and come. On the second he asks the Lord to stir up our
hearts. A prayer we know from the Fourth Sunday appears there after the winter
Ember Days (\emph{Gregorian Sacramentary}, nos. 150--157). In the ninth
century Amalarius of Metz knew older books with five weeks and a book he
believed to be Gregory the Great's with four (\emph{On the Liturgy} 3.40).

In the 1962 books Advent begins with First Vespers of the Sunday nearest
November 30. Vespers is the evening prayer of the Church's Office. First
Vespers is prayed on the eve of a Sunday or feast. Advent ends after None, the
afternoon Hour, on December 24. The Church begins to keep Christmas at its
First Vespers that evening. The Missal has four Sunday Masses for Advent. If
the Vigil of the Nativity falls on a Sunday, its Mass takes the place of the
Fourth Sunday (\emph{General Rubrics}, nos. 20, 30, 71--72). The Sunday Masses
of Advent ordinarily use violet. They omit the Gloria, the hymn that begins
with the angels' words at Bethlehem, and still sing the Alleluia.
